\section{What the Gaussian Assumption Proves}
%% ============================================================

Suppose the public signal is drawn from $Y_t\sim\mathcal{N}(\mu_C,\sigma^2)$ under cooperation and $Y_t\sim\mathcal{N}(\mu_D,\sigma^2)$ after the specified deviation, and punishment is triggered when $Y_t<\tau$ for a fixed threshold $\mu_D<\tau<\mu_C$. Then
\[
\alpha = \Phi\!\left(\frac{\tau-\mu_C}{\sigma}\right),\qquad
\beta  = \Phi\!\left(\frac{\tau-\mu_D}{\sigma}\right).
\]

\begin{proposition}[Restricted noise comparative static]\label{prop:gauss}
Holding $\mu_C$, $\mu_D$, $\tau$, prescribed actions, and payoff gaps fixed with $\mu_D<\tau<\mu_C$, increasing noise $\sigma>0$ strictly decreases $\Delta=\beta-\alpha$ and strictly raises $\delta^*$ in~\eqref{eq:threshold}.
\end{proposition}

\begin{proof}
Let $z_C=(\tau-\mu_C)/\sigma < 0$ and $z_D=(\tau-\mu_D)/\sigma > 0$. With standard normal density $\phi$,
\begin{equation}\label{eq:ddelta}
\frac{\partial\Delta}{\partial\sigma}
= -\frac{1}{\sigma}\bigl[z_D\,\phi(z_D)-z_C\,\phi(z_C)\bigr] < 0,
\end{equation}
since $\alpha'(\sigma)=-z_C\phi(z_C)/\sigma > 0$ and $\beta'(\sigma)=-z_D\phi(z_D)/\sigma < 0$. Combining~\eqref{eq:ddelta} with~\eqref{eq:compstat} yields $\partial\delta^*/\partial\sigma = g[(g+\Delta V)\alpha'(\sigma)-\Delta V\,\beta'(\sigma)]/D^2 > 0$. As $\sigma\to 0^+$, $\alpha\to 0$ and $\beta\to 1$, so $\delta^*\to g/(g+\Delta V)$. As $\sigma\to\infty$, both $\alpha,\beta\to 1/2$, so $\delta^*\to 2$, which is infeasible for $\delta<1$ (Appendix~A.3). \qed
\end{proof}

\subsection{Conditional Connection to the PPO Environment}

Propositions~\ref{prop:threshold} and~\ref{prop:gauss} provide an exact incentive accounting benchmark, but three structural differences separate this analytical model from the PPO experiment: (i) the PPO simulator delivers continuous, two-period-averaged, one-step-delayed reports to a 32-step GRU rather than enforcing a contemporaneous binary threshold; (ii) the experiment compares an informative channel against an action-independent replacement rather than sweeping $\sigma$ across a fixed trigger policy; and (iii) PPO episodes terminate at $T=512$ with observable normalized time rather than running over an infinite horizon. Consequently, our empirical audit tests whether trained recurrent policies exhibit continuation-value deterrence in practice, without claiming that finite-horizon PPO must replicate infinite-horizon trigger equilibria.

%% ============================================================
